% Preámbulo adaptado del template. Compilador: XeLaTeX.
\usepackage{fontspec}
\usepackage[spanish,es-noquoting,es-nodecimaldot]{babel} % Idioma de la clase

% Paquetes matemáticos y gráficos
\usepackage{amsmath, amssymb}
\usepackage{graphicx}
\graphicspath{{images/}}
\usepackage{geometry}
\usepackage{multicol}

% Paquetes para el tema personalizado (del código proporcionado)
\usepackage{tikz,etoolbox}
\usetikzlibrary{shapes,arrows,arrows.meta,positioning,calc,fit,backgrounds}
\usepackage{booktabs,array}
\usepackage{listings}

% Configuración de referencias bibliográficas
\usepackage[style=apa, backend=biber]{biblatex}
\addbibresource{4_bibliography.bib}
\DeclareLanguageMapping{spanish}{spanish-apa}
\usepackage{csquotes}
\usepackage{pgfplots}
\pgfplotsset{compat=1.18}

% Paquete para que las páginas de apéndices no cuenten en la numeración
\usepackage{appendixnumberbeamer}

% Definir colores personalizados (del código proporcionado)
\definecolor{UWBlack}{HTML}{000000}
\definecolor{UWWhite}{HTML}{FFFFFF}
\definecolor{UWMathPinkL1}{HTML}{FFBEEF}
\definecolor{UWMathPinkL2}{HTML}{FF63AA}
\definecolor{UWMathPinkL3}{HTML}{DF2498}
\definecolor{UWMathPinkL4}{HTML}{C60078}
\definecolor{UWGrayL1}{HTML}{DFDFDF}
\definecolor{UWGrayL2}{HTML}{A2A2A2}
\definecolor{UWGrayL3}{HTML}{787878}
\definecolor{UWGrayL4}{HTML}{000000}
\definecolor{UWGoldL1}{HTML}{FFFFAA}
\definecolor{UWGoldL2}{HTML}{FFEA3D}
\definecolor{UWGoldL3}{HTML}{FFD54F}
\definecolor{UWGoldL4}{HTML}{E4B429}
\definecolor{carrot}{HTML}{EE693F}
\definecolor{ivory}{HTML}{F1F3CE}
\definecolor{emerald}{HTML}{265C00}
\definecolor{turquise}{HTML}{5BC8AC}
\definecolor{peacockblue}{HTML}{1E656D}
\definecolor{spicy}{HTML}{B51D0A}
\definecolor{bluegreen}{HTML}{5F968E}
\definecolor{rust}{HTML}{9B4F0F}
\definecolor{burntorange}{HTML}{DE7A22}
\definecolor{sea}{HTML}{20948B}
\definecolor{lagoon}{HTML}{6AB187}

% Colores para los entornos bluebox, greenbox, redbox (definidos previamente)
\definecolor{myblue}{RGB}{173,216,230}
\definecolor{mygreen}{RGB}{144,238,144}
\definecolor{myred}{RGB}{255,182,193}

% Configuración de colores para los componentes de las diapositivas (del código proporcionado)
\setbeamercolor{background canvas}{bg=UWWhite}
\setbeamercolor{author}{fg=UWGoldL3}
\setbeamercolor{institute}{fg=UWMathPinkL3}
\setbeamercolor{title}{fg=UWGrayL4}
\setbeamercolor{section in head/foot}{bg=UWBlack, fg=UWGoldL3}
\setbeamercolor{author in head/foot}{fg=UWGoldL3, bg=UWBlack}
\setbeamercolor{title in head/foot}{fg=UWBlack, bg=UWGoldL3}
\setbeamercolor{institute in head/foot}{fg=UWGoldL3, bg=UWBlack}
\setbeamercolor{navigation symbols}{fg=UWBlack}
\setbeamercolor{normal text}{fg=UWGrayL4}
\setbeamercolor{section in toc}{fg=emerald}
\setbeamercolor{subsection in toc}{fg=bluegreen}
\setbeamercolor{frametitle}{fg=UWMathPinkL4, bg=UWGrayL1}
\setbeamercolor{block title}{bg=emerald, fg=ivory}
\setbeamercolor{block body}{bg=peacockblue!20, fg=peacockblue}
\setbeamercolor{section number projected}{bg=turquise, fg=black}
\setbeamercolor{block title example}{fg=rust, bg=sea!40}
\setbeamercolor{block body example}{fg=burntorange, bg=lagoon!20}

% Configuración de fuentes
\setbeamerfont{frametitle}{series=\bfseries} % Título de diapositiva en negrita
\setbeamerfont{section number projected}{family=\rmfamily, series=\bfseries, size=\normalsize} % Números de sección en TOC

% Contadores para Definición, Teorema y Ejemplo
\newcounter{defcounter}
\newcounter{theocounter}
\newcounter{excounter}

% Redefinir los entornos bluebox, greenbox y redbox con contadores automáticos
\newenvironment{bluebox}[1]
{%
  \stepcounter{defcounter}% Incrementar el contador de definiciones
  \setbeamercolor{block title}{bg=myblue, fg=black}%
  \begin{block}{Definición \thedefcounter\ (#1)}%
}
{\end{block}}

\newenvironment{greenbox}[1]
{%
  \stepcounter{theocounter}% Incrementar el contador de teoremas
  \setbeamercolor{block title}{bg=mygreen, fg=black}%
  \begin{block}{Teorema \thetheocounter\ (#1)}%
}
{\end{block}}

\newenvironment{redbox}[1]
{%
  \stepcounter{excounter}% Incrementar el contador de ejemplos
  \setbeamercolor{block title}{bg=myred, fg=black}%
  \begin{block}{Ejemplo \theexcounter\ (#1)}%
}
{\end{block}}

% Comandos para campos comunes en presentaciones (del código proporcionado)
\newcommand\course[1]{\def\insertcourse{#1}}
\newcommand\class[1]{\def\insertclass{#1}}



% Configuración de la diapositiva de título (del código proporcionado)
\setbeamertemplate{title page}{
\begin{tikzpicture}[remember picture,overlay]
  \fill[UWWhite] ([yshift=30pt]current page.west) rectangle (current page.south east);
  \fill[UWBlack] ([yshift=30pt]current page.west) rectangle (current page.north east);
  \node[anchor=east] at ([yshift=-52pt,xshift=-15pt]current page.north east)
    {\includegraphics[width=0.34\paperwidth]{UP_logo.jpg}};
  \node[anchor=north west] at ([yshift=-35pt,xshift=15pt]current page.north west)
    {\parbox[t]{0.51\paperwidth}{\usebeamercolor[fg]{author}\fontsize{\TamNombrePortada}{17}\selectfont\bfseries\insertauthor\textsuperscript{*}}};
  \node[anchor=north west] at ([yshift=-61pt,xshift=15pt]current page.north west)
    {\parbox[t]{0.51\paperwidth}{\hypersetup{urlcolor=UWWhite}\small\href{mailto:al.venturam@up.edu.pe}{al.venturam@up.edu.pe}}};
  \node[anchor=north west] at ([yshift=-82pt,xshift=15pt]current page.north west)
    {\parbox[t]{0.51\paperwidth}{\usebeamercolor[fg]{institute}\normalsize\insertinstitute}};
  \node[anchor=north] at ([yshift=-1pt]current page.center)
    {\parbox[t]{\textwidth}{\centering\usebeamercolor[fg]{title}\fontsize{\TamTituloPortada}{\InterlineadoTituloPortada}\selectfont\bfseries\inserttitle}};
  \node[anchor=north] at ([yshift=-60pt]current page.center)
    {\parbox[t]{\textwidth}{\centering\fontsize{\TamSubtituloPortada}{15}\selectfont\SubtituloPortada}};
  \node[anchor=north] at ([yshift=-83pt]current page.center)
    {\parbox[t]{\textwidth}{\centering\fontsize{\TamCursoPortada}{14}\selectfont\insertcourse}};
  \node[anchor=north] at ([yshift=-100pt]current page.center)
    {\parbox[t]{\textwidth}{\centering\fontsize{\TamSemestrePortada}{12}\selectfont\insertclass}};
  \node[anchor=south] at ([yshift=10pt]current page.south)
    {\parbox{0.92\paperwidth}{\centering\fontsize{7}{9}\selectfont\textsuperscript{*}\NotaExpositor}};
\end{tikzpicture}
}

% Configuración de la cabecera (del código proporcionado)
\setbeamertemplate{headline}
{%
  \begin{beamercolorbox}{section in head/foot}
    \vskip5pt\bfseries
    \insertnavigation{\paperwidth}
    \vskip2pt
  \end{beamercolorbox}%
}

\renewcommand*{\slideentry}[6]{} % Sin círculo sólido en la cabecera

% Configuración del pie de página (del código proporcionado)
\setbeamertemplate{footline}
{
	\leavevmode
	\begin{beamercolorbox}[wd=.333333\paperwidth,ht=2.5ex,dp=1.125ex,
      leftskip=.3cm,rightskip=.3cm plus1fil]{author in head/foot}
		\usebeamerfont{author in head/foot}\insertshortauthor
    \end{beamercolorbox}%
    \begin{beamercolorbox}[wd=.333333\paperwidth,ht=2.5ex,dp=1.125ex,
      leftskip=.3cm,rightskip=.3cm plus1fil,center]{title in head/foot}
      {\usebeamerfont{title in head/foot}\insertshorttitle}
    \end{beamercolorbox}%
    \begin{beamercolorbox}[wd=.333333\paperwidth,ht=2.5ex,dp=1.125ex,
      leftskip=.3cm,rightskip=.3cm plus1fil]{institute in head/foot}
      \hfill    {\usebeamercolor[fg]{institute in head/foot}\insertshortinstitute}    
	\end{beamercolorbox}%
}

% Configuración de los símbolos de navegación (número de diapositiva)
\setbeamertemplate{navigation symbols}{\bfseries\insertframenumber/\inserttotalframenumber}

% Estilo de las secciones/subsecciones en la tabla de contenidos
\setbeamertemplate{sections/subsections in toc}[ball]

% Estilo de las viñetas de itemize (del código proporcionado)
\setbeamertemplate{itemize item}{
	\tikz{
		\draw[fill=spicy,draw=none] (0, 0) rectangle(0.075, 0.075);
		\draw[fill=spicy,draw=none] (0.075, 0.075) rectangle(0.15, 0.15);
		\draw[fill=spicy,draw=none] (0, 0.15) rectangle(0.075, 0.225);
	}
}

\setbeamertemplate{itemize subitem}{
	\tikz{
		\draw[fill=carrot,draw=none] (0, 0) rectangle(0.05, 0.05);
		\draw[fill=carrot,draw=none] (0.05, 0.05) rectangle(0.1, 0.1);
		\draw[fill=carrot,draw=none] (0, 0.1) rectangle(0.05, 0.15);
	}
}

\AtBeginEnvironment{block}{
	\setbeamertemplate{itemize item}{
		\tikz{
			\draw[fill=spicy,draw=none] (0, 0) rectangle(0.075, 0.075);
			\draw[fill=spicy,draw=none] (0.075, 0.075) rectangle(0.15, 0.15);
			\draw[fill=spicy,draw=none] (0, 0.15) rectangle(0.075, 0.225);
		}
	}

	\setbeamertemplate{itemize subitem}{
		\tikz{
			\draw[fill=carrot,draw=none] (0, 0) rectangle(0.05, 0.05);
			\draw[fill=carrot,draw=none] (0.05, 0.05) rectangle(0.1, 0.1);
			\draw[fill=carrot,draw=none] (0, 0.1) rectangle(0.05, 0.15);
		}
	}
}

% Configuración de bloques
\setbeamertemplate{blocks}[rounded][shadow=false]

% Configuración de transparencia para elementos cubiertos
\setbeamercovered{invisible}

% Cambiar la fuente de matemáticas a serif


% Configuración de notas al pie
\renewcommand{\thefootnote}{\ensuremath{\ast}}

% Configuración de la diapositiva de título
\title[Crash course de Dynare]{Un crash course\\de Dynare}
\author[Alejandro V.]{Alejandro Ventura}
\institute[UP]{Universidad del Pacífico}
\course{Macroeconomía II}
\class{2026 II}
\setbeameroption{hide notes}
\setbeamertemplate{bibliography item}{}
\setbeamercolor{bibliography entry author}{fg=UWBlack}
\setbeamercolor{bibliography entry title}{fg=UWBlack}
\setbeamercolor{bibliography entry location}{fg=UWBlack}
\setbeamercolor{bibliography entry note}{fg=UWBlack}
\hypersetup{colorlinks=true,linkcolor=UWBlack,citecolor=UWBlack,urlcolor=peacockblue}
\setbeamercolor{description item}{fg=UWMathPinkL4}
\setbeamercolor{item projected}{fg=UWMathPinkL4}
\newcommand{\fuente}[1]{\par\vspace{0.15cm}{\scriptsize #1\par}}
% EDITA AQUÍ LA TIPOGRAFÍA Y LOS TAMAÑOS DE PORTADA.
% La misma familia se usa en texto, títulos y citas.
\newcommand{\FuenteTexto}{Latin Modern Sans}
\newcommand{\TamTituloPortada}{22}
\newcommand{\InterlineadoTituloPortada}{25}
\newcommand{\TamNombrePortada}{14}
\newcommand{\TamSubtituloPortada}{12}
\newcommand{\TamCursoPortada}{11}
\newcommand{\TamSemestrePortada}{10}
\newcommand{\NotaExpositor}{Todo lo dicho en esta sesión es responsabilidad del expositor.}
\newcommand{\SubtituloPortada}{Del papel al código: el modelo RBC}

\setmainfont{\FuenteTexto}
\setsansfont{\FuenteTexto}
\renewcommand{\familydefault}{\sfdefault}
\usefonttheme{professionalfonts}
\setbeamerfont{normal text}{family=\sffamily,size=\normalsize}
\setbeamerfont{frametitle}{family=\sffamily,series=\bfseries}
\setbeamerfont{quote}{family=\sffamily,shape=\upshape,size=\normalsize}
\setbeamerfont{quotation}{family=\sffamily,shape=\upshape,size=\normalsize}

% Iconos de una biblioteca existente; no representan hechos históricos.
\usepackage{fontawesome5}

% Todos los rótulos son visibles: no heredan el color negro de los enlaces.
\setbeamertemplate{headline}{%
  \begin{beamercolorbox}[wd=\paperwidth,ht=3ex,dp=1ex]{section in head/foot}
    \sffamily\bfseries\hbox to\paperwidth{\hfil Motivación\hfil Modelo\hfil Dynare\hfil Resultados\hfil Cierre\hfil}
  \end{beamercolorbox}
}


% ---- Añadidos para el crash course de Dynare --------------------------------
% Colores del handout, para que las figuras se vean iguales en ambos materiales.
\definecolor{darkblue}{rgb}{0.0, 0.0, 0.55}
\definecolor{thmblue}{RGB}{228, 238, 255}
\definecolor{upblue}{RGB}{15, 35, 85}
\definecolor{errframe}{RGB}{193, 42, 32}
\definecolor{tipframe}{RGB}{0, 128, 102}
\definecolor{codebg}{RGB}{246, 246, 246}
% Código de Dynare
\lstdefinelanguage{dynare}{
  morekeywords={var,varexo,parameters,model,end,steady_state_model,shocks,stderr,
    steady,resid,check,stoch_simul,initval,linear,log,exp},
  sensitive=true, morecomment=[l]{//}, morestring=[b]'}
\lstdefinestyle{dyn}{language=dynare, basicstyle=\ttfamily\footnotesize,
  keywordstyle=\color{darkblue}\bfseries, commentstyle=\color{UWGrayL3}\itshape,
  stringstyle=\color{rust}, backgroundcolor=\color{codebg}, frame=single,
  rulecolor=\color{UWGrayL1}, columns=fullflexible, keepspaces=true,
  xleftmargin=2pt, xrightmargin=2pt, aboveskip=4pt, belowskip=4pt,
  showstringspaces=false, literate={á}{{\'a}}1 {é}{{\'e}}1 {í}{{\'i}}1 {ó}{{\'o}}1 {ú}{{\'u}}1 {ñ}{{\~n}}1}
\lstset{style=dyn}
\newcommand{\cod}[1]{{\ttfamily\color{darkblue}#1}}
% Cifra destacada
\newcommand{\num}[1]{\textcolor{UWMathPinkL4}{\textbf{#1}}}
